\section{Conclusion}
\label{sec:conclusion}

Provenance graph evasion attacks face a dual failure problem: injected operations may corrupt the attack's causal dependencies (\emph{safety failure}) or introduce statistical anomalies that make the evasion itself detectable (\emph{stealth failure}).
We presented \textsc{SafeMimic}, a framework that addresses both failure modes through a principled safety-first design.
By formally defining a dynamic safety space---the set of operations whose system effects are provably disjoint from the attack's dependency set---\textsc{SafeMimic} guarantees that injected camouflage never disrupts attack functionality (Safety Theorem).
Within this safe space, distribution-aware selection via KL-divergence minimization ensures that the resulting provenance graph is statistically indistinguishable from benign activity.

Our evaluation demonstrates that \textsc{SafeMimic} achieves 100\% attack preservation and 0\% ML detection across both feature-based and GNN-based detectors.
In contrast, we show that ProvNinja's regularity-based approach---which lacks safety reasoning---produces a \emph{backfire effect}: its injected operations create frequency anomalies and dependency violations that increase rather than decrease detectability.
We formalize this finding as a proposition establishing that regularity does not imply safety, providing theoretical grounding for an empirically observed but previously unexplained phenomenon.

The key contributions of this work are: (1)~the first formal treatment of safety in provenance graph evasion, (2)~the dynamic safety space abstraction that adapts as operations are injected, (3)~theoretical results connecting safety, regularity, and detectability, and (4)~empirical evidence that safety-unaware evasion can backfire against modern detectors.

Several directions remain for future work.
First, validating \textsc{SafeMimic} on the full-scale authentic DARPA TC datasets (rather than calibrated synthetic graphs) would strengthen the empirical findings.
Second, investigating adaptive defenses that are specifically aware of distribution-matching evasion strategies---for instance, detectors that monitor for unnaturally smooth operation distributions---is an important next step for the arms race.
Third, extending the framework to additional operating systems (e.g., Windows ETW, macOS DTrace) and cloud VM infrastructures would demonstrate the generality of the safety-space abstraction in practice.
